\documentclass[aspectratio=169]{beamer}
\usetheme{Boadilla}
\setbeamertemplate{navigation symbols}{}
\usepackage{tikz}
\usepackage{amsmath}
\usepackage{stmaryrd}
\usepackage{mathtools}

\title[Correct-by-Construction $\mu$-Calculus]{Rational Codata as Syntax-with-Binding}
\subtitle{Correct-by-Construction Foundations of the Modal $\mu$-Calculus}
\author[S. Watters]{Sean Watters}
\institute{University of Strathclyde}
\date{9th June 2025}


\newcommand{\meaningof}[1]{\ensuremath{\llbracket #1 \rrbracket}}
\newcommand{\hiddencell}[2]{\color{red} \action<#1->{#2}}

\begin{document}


\begin{frame}
  \titlepage{}
\end{frame}

\begin{frame}{Plan}
  We will:
\begin{enumerate}
\item Introduce the $\mu$-calculus, and its Fischer-Ladner closure.
        \smallskip
\item Sketch our (now complete!) formalised proof of the closure's finiteness.
        \smallskip
\item Discuss the presentation of rational cotrees as syntax with binding, its role in the proof, and future plans in this direction.
        \smallskip
\item Aim to keep it high-level and skip the gory details!
\end{enumerate}


\end{frame}

% \begin{frame}{Kripke Semantics}

% \begin{definition}
%   A \textbf{Kripke model} is a tuple $(S,T,V)$ of a set of states $S$, a transition relation $T : S \to S \to \mathsf{Prop}$, and a valuation function $V : A \to S \to \mathsf{Prop}$ (given some set $A$ of atomic propositions).
% \end{definition}

% \bigskip

% Notice that:
% \begin{itemize}
%   \item $S$ and $T$ induce a (possibly infinite) graph-like structure.
%         \smallskip
%   \item Each state $s \in S$ is equipped with a model of propositional logic, $a \mapsto V~a~s$.
%         \smallskip
%   \item Partial application of $T$ on some state $s \in S$ yields the ``is a successor of $s$'' predicate.
% \end{itemize}

% \end{frame}

% \begin{frame}{Modal Logic}
%   \textbf{Syntax:}
%   For all propositional atoms $a \in \mathsf{At}$:
%   \begin{equation*}
%     \varphi := a~|~\neg \varphi~|~\varphi \land \varphi~|~\varphi \lor \varphi~|~\varphi \Rightarrow \varphi~|~\Box \varphi~|~\Diamond \varphi
%   \end{equation*}

%   \bigskip

%   \textbf{Semantics:}
%   Given a Kripke model $\mathcal{M} = (S,T,V)$, let $\meaningof{-}~- : ML \to S \to \mathsf{Prop}$, where:
%   \begin{align*}
%     \meaningof{a}~s &:= V~a~s \\
%     \meaningof{\lnot \varphi}~s &:= \lnot~(\meaningof{\varphi}~s) \\
%     \meaningof{\varphi \land \psi}~s &:= (\meaningof{\varphi}~s) \times (\meaningof{\psi}~s) \\
%     \meaningof{\varphi \lor \psi}~s &:= (\meaningof{\varphi}~s) + (\meaningof{\psi}~s) \\
%     \meaningof{\varphi \Rightarrow \psi}~s &:= (\meaningof{\varphi}~s) \to (\meaningof{\psi}~s) \\
%     \meaningof{\Box \varphi}~s &:= \forall t \in (T~s).~\meaningof{\varphi}~t \\
%     \meaningof{\Diamond \varphi}~s &:= \exists t \in (T~s).~\meaningof{\varphi}~t
%   \end{align*}
% \end{frame}


\begin{frame}{The Modal $\mu$-Calculus}
  \textbf{Syntax:}
  For all propositional atoms $a \in \mathsf{At}$ and variable names $x \in \mathsf{Var}$:
  \begin{equation*}
    \varphi := a~|~\neg a~|~x~|~\varphi \land \varphi~|~\varphi \lor \varphi~|~\Box \varphi~|~\Diamond \varphi~|~\mu x. \varphi~|~\nu x. \varphi
  \end{equation*}

  \bigskip

  Notes:
 \begin{itemize}
  \item The fixpoint operators $\mu$ and $\nu$ are variable binders.
         \smallskip
  \item The syntax is strictly positive --- this matters when giving semantics to fixpoint formulas.
         \smallskip

 \end{itemize}
 \bigskip
 \textbf{Semantics:}
 \begin{itemize}
  \item Kripke semantics, where $\mu$ and $\nu$ let us reason about unbounded/infinite behaviour.
         \smallskip
  \item Satisfiability and model checking are decidable.
         \smallskip
  \item The $\mu$-calculus subsumes temporal and dynamic logics such as LTL, CTL*, and PDL.
 \end{itemize}

  % \textbf{Semantics}:
  % Given a Kripke model $\mathcal{M}$ and an interpretation of free variables $i : \mathsf{Var} \to S \to \mathsf{Prop}$, let:
  % \begin{align*}
  %   \meaningof{x}_{i}~s &:= i~x~s \\
  %   \meaningof{\mu x.~\varphi}_{i}~s &:= \mu~((U : S \to \mathsf{Prop}) \mapsto \meaningof{\varphi}_{i[x:=U]})~s \\
  %   \meaningof{\nu x.~\varphi}_{i}~s &:= \nu~((U : S \to \mathsf{Prop}) \mapsto \meaningof{\varphi}_{i[x:=U]})~s
  % \end{align*}

\end{frame}

\begin{frame}{Fixpoint Unfolding}
  At the heart of the $\mu$-calculus is the semantic equivalence:
    \begin{equation*}
      \eta x.~\varphi~\equiv~\varphi [ x := \eta x.\varphi]
    \end{equation*}

  We call $\varphi [ x := \eta x.\varphi ]$ the \textbf{unfolding} of $\eta x.~\varphi$.

\bigskip

  For example: let $E(p) := \mu x.~p \lor \Diamond x$. Then:
  \begin{align*}
    &E(p) \\
    \equiv~&p \lor \Diamond(E(p)) \\
    \equiv~&p \lor \Diamond(p \lor \Diamond(E(p))) \\
    \equiv~&\ldots
  \end{align*}

  \bigskip

  % This being a \emph{least} fixpoint says that the formula will be validated in \emph{finitely} many unfoldings.

  % it's the whole point
  % substitution (hand wavy, will do properly with dB in part 2)
  % example showing formula semantically equivalent to its unfolding. Show repeated unfoldings to intuit mu/nu
\end{frame}

\begin{frame}{The Closure}
  \begin{definition}
    The \textbf{closure} of a formula $\varphi$ is the minimal set which contains $\varphi$, and is closed under taking unfoldings of fixpoint formulas, and direct subformulas of non-fixpoint formulas.
  \end{definition}

  \bigskip

  In other words, it is the minimal set $C$ satisfying:
    \begin{align*}
      &\varphi \in C \\
      \bigcirc \varphi \in C \Rightarrow~&\varphi \in C,~\mathsf{where}~\bigcirc \in \{\Box,\Diamond\} \\
      \varphi \star \psi \in C \Rightarrow~&\varphi \in C~\mathsf{and}~\psi \in C,~\mathsf{where}~\star \in \{\land,\lor\} \\
      \eta x. \varphi \in C \Rightarrow~&\varphi [ x := \mu x. \varphi] \in C,~\mathsf{where}~\eta \in \{\mu,\nu\}
    \end{align*}
  % insability under alpha equivalence
  % relationship to model checking
\end{frame}

\begin{frame}{Finiteness of the Closure (1)}
  \begin{theorem}
    For all $\varphi$, the closure of $\varphi$ is finite. (Kozen, 1983)
  \end{theorem}

  \bigskip

  \textbf{Proof Sketch (Kozen):}
  \smallskip

 \begin{enumerate}
   \item Define the \emph{expansion} of a formula as the sequential instantiation of all its free variables.
         \smallskip
   \item Define an alternative, structurally inductive procedure for computing the closure via the expansion map.
         \smallskip
   \item Prove the alternative definition correct by induction.
 \end{enumerate}

 \bigskip

 \begin{center}
 \textbf{Not so simple in a formal setting!}
 \end{center}

\end{frame}

\begin{frame}{Finiteness of the Closure (2)}
  \textbf{Proof Sketch (Our Approach in Agda):}
  \smallskip

 \begin{enumerate}
   \item Implement the closure coinductively as a nonwellfounded cotree, via the standard definition. (Trivially correct).
         \smallskip
   \item Implement the alternative definition of the closure as an inductive syntax-with-binding representation of a rational cotree (a ``tree with backedges''). (Finite by construction).
         \smallskip
   \item Define the ``unfolding'' of such a tree-with-backedges to a cotree.
         \smallskip
   \item Prove the coinductive definition closure bisimilar to the unfolding of the inductive approximation.
         \smallskip
   \item Prove that we can transport the correctness proofs across the bisimulation, to show that the finite-by-construction algorithm does in fact compute the closure.
 \end{enumerate}
\end{frame}

\begin{frame}{Nonwellfounded Cotrees}
  % define trees in agda
\end{frame}

\begin{frame}{Rational Trees}
  % cite the loopy bois
  % define scopes and trees in agda
  % mention equivalence up to renaming, and that this + a thinning between the scopes gets you structural equivalence of trees.
\end{frame}

\begin{frame}{Unfolding Trees}
  % define unfolding in agda
\end{frame}

\begin{frame}{Progress and Future Work}
\begin{itemize}
        \item The proof is now finished, and the Agda formalisation is on my GitHub.
        \smallskip
        \item What's the general, abstract story about this presentation of rational codata? (Rational fixpoint of containers??)
        \smallskip
        \item IMO: The $\mu$-calculus and type theory are ripe for more cross-fertilisation.

\end{itemize}

\bigskip

\begin{center}
  \Large Thanks!
\end{center}
\end{frame}

\end{document}
